%++++++% Preâmbulo %+++++++++++++++++++++++++++++++++++++++++++++++++++++++++
\documentclass[13pt, xcolor={dvipsnames,svgnames}, portuguese]{beamer}
%\documentclass[[11pt, xcolor={dvipsnames,svgnames,table},portuguese]{beamer} 

\usetheme{CambridgeUS}

\setbeamercolor*{structure}{bg=PineGreen!20,fg=PineGreen} %fg=PineGreen
\definecolor{beamer@pinegreen}{rgb}{0.137,0.666,0.741}


\setbeamercolor*{palette primary}{use=structure,fg=white,bg=structure.fg}
\setbeamercolor*{palette secondary}{use=structure,fg=white,bg=structure.fg!75}
\setbeamercolor*{palette tertiary}{use=structure,fg=white,bg=structure.fg!50!black}
\setbeamercolor*{palette quaternary}{fg=white,bg=black}

\setbeamercolor{section in toc}{fg=black,bg=white}
\setbeamercolor{alerted text}{use=structure,fg=structure.fg!50!black!80!black}

\setbeamercolor{titlelike}{parent=palette primary,fg=structure.fg!50!black}
\setbeamercolor{frametitle}{bg=gray!10!white,fg=PineGreen}

\setbeamercolor*{titlelike}{parent=palette primary}

\usepackage[utf8]{inputenc}
\usepackage[brazil]{babel}  % idioma
\usepackage{amsmath,amsfonts,amssymb,textcomp}
\usepackage{graphicx}
\usepackage{subfigure}
\usepackage[utf8]{inputenc}
\usepackage{ifpdf}
\usepackage{listings}
\usepackage{dirtree}

% Configurações para o ambiente lstlisting
\lstset{
    language=C,
    basicstyle=\ttfamily\footnotesize,
    numbers=left,
    numberstyle=\tiny,
    numbersep=5pt
}
\lstdefinelanguage{JavaScript}{
  keywords={typeof, new, true, false, catch, function, return, null, catch, switch, var, if, in, while, do, else, case, break},
  keywordstyle=\color{blue}\bfseries,
  ndkeywords={class, export, boolean, throw, implements, import, this},
  ndkeywordstyle=\color{darkgray}\bfseries,
  identifierstyle=\color{black},
  sensitive=false,
  comment=[l]{//},
  morecomment=[s]{/*}{*/},
  commentstyle=\color{purple}\ttfamily,
  stringstyle=\color{red}\ttfamily,
  morestring=[b]',
  morestring=[b]"
}
% here you should include other packages with \usepackage

    \ifpdf

      % hyperref should be the last package loaded:
	  %\usepackage[pdftex]{hyperref}
      \usepackage{pst-pdf}
    \else

      % make the command \href from hyperref available as a 'print only'
      \newcommand{\href}[2]{#2}

    \fi

%Global Background must be put in preamble
\usebackgroundtemplate%
{%
    \includegraphics[width=\paperwidth,height=\paperheight]{Figuras/fundo.png}%
}
\setbeamertemplate{frametitle}[default][center]
 
\author{Othon Oliveira}
\title{Introdução ao React.JS}
\institute{SENAC - PROA} 
\date{} 
%\subject{} 

\begin{document}

\begin{frame}
\titlepage
%\date{}
\end{frame}

% Capa - requer o TikZ
\newcommand{\capa}{
    \begin{tikzpicture}[remember picture,overlay]
        \node at (current page.south west)
            {\begin{tikzpicture}[remember picture, overlay]
                \fill[shading=radial,top color=orange,bottom color=orange,middle color=yellow] (0,0) rectangle (\paperwidth,\paperheight);
            \end{tikzpicture}
          };
    \end{tikzpicture}
}


%+++++++++++++++++++++++ Telas iniciais  ---------------------------------
\begin{frame}\frametitle{Sumário}
\tableofcontents
\end{frame}
%---------------------------------------------------------------------
\begin{frame}
\frametitle{Minha Imagem}
\begin{figure}
\centering
\includegraphics[width=0.5\textwidth]{Figuras/react1.png}
\caption{Componentes React}
\end{figure}
\end{frame}
\begin{frame}\frametitle{Sumário}
\tableofcontents
\end{frame}


%+++++++++++++++++++++++++++++++++++++++++++++++
\section{Entrada de dados}
%+++++++++++++++++++++++++++++++++++++++++++++++
\begin{frame}{Criar a seguinte árvore de diretórios}
\frametitle{Estrutura de Diretórios do Projeto - Exemplo}
\framesubtitle{Revendo hooks}

\dirtree{%
.1 seu-app/.
.2 src/.
.3 App.js.
.3 index.js.
.3 logo.svg.
.3 paginas/.
.4 Home.js.
.4 Sobre.js.
.3 componentes/.
.4 HookState.js.
.2 .gitignore.
}
\end{frame}


%+++++++++++++++++++++++++++++++++++++++++++++++
\section{Funções Especiais}
%+++++++++++++++++++++++++++++++++++++++++++++++
\begin{frame}
\frametitle{Explicação do \texttt{useState}}
Neste exemplo, vamos explorar como usar o hook \texttt{useState} para gerenciar estados em componentes funcionais do React. Vamos criar um contador simples que aumenta sempre que o botão é clicado.

Bora codar?
\end{frame}

%-------------------------------------------------------------


\begin{frame}
\frametitle{Explicação do  \texttt{useEffect}}
Neste exemplo, aprenderemos como usar o hook \texttt{useEffect} para buscar renderizar uma parte da página. 

Bora codar?
\end{frame}

%+++++++++++++++++++++++++++++++++++++++++++++++
\section{Calculos com }
%+++++++++++++++++++++++++++++++++++++++++++++++

%-------------------------------------------------------------


%+++++++++++++++++++++++++++++++++++++++++++++++
\section{Calculos com }
%+++++++++++++++++++++++++++++++++++++++++++++++

%-------------------------------------------------------------

%-------------------------------------------------------------


%%+++++++++++++++++++++++++++++++++++++++++++++++


%+++++++++++++++++++++++++++++++++++++++++++++++




%-----------------------------------------------

\end{document}
